% META: {"id": "1NSI-RESEAUX-RE-C3", "chapitre": "1NSI-RESEAUX", "type_objet": "remediation", "capacites": ["P-ARCH-02C"], "capacites_codes": ["C3"], "mode_creation": "ex_nihilo", "sources_inspiration": [], "status": "needs_review"}

%---------------------------------------------------------------
% FICHE DE REMEDIATION — C3 : un cable se note des deux cotes
%---------------------------------------------------------------

\textbf{Ce qui bloque.} Ton dictionnaire dit qu'une machine « voit » sa voisine, mais pas
l'inverse : la recherche de chemin répond alors \lstinline{True} dans un sens et
\lstinline{False} dans l'autre. Le réseau n'a pas changé, c'est le modèle qui est faux. Et
quand un cycle apparaît, tu crois devoir l'interdire alors que le marquage des visités
suffit.

\textbf{À retenir.} Un câble relie deux machines : il figure dans la liste de \emph{chacune}.
Toute machine citée comme voisine doit aussi être une clé. La recherche s'arrête toujours,
cycle ou non, parce qu'une machine visitée n'est jamais réexplorée.

\begin{exercice}{1NSI-RES-RE-C3-EX1}{1}{12}
Un élève modélise le réseau de sa maison : une box reliée à un PC, à une télévision et à une
tablette.
% PYTHON-SOURCE: code/1NSI-RESEAUX-RE-C3.py
\begin{python}
reseau = {"Box": ["PC", "TV", "Tablette"], "PC": ["Box"], "TV": []}
\end{python}
\begin{enumerate}
  \item Que renvoie \lstinline{peuvent_communiquer(reseau, "Box", "TV")} ? Et
        \lstinline{peuvent_communiquer(reseau, "TV", "Box")} ? Expliquer la contradiction.
  \item Corriger le dictionnaire : quelles listes compléter, quelle clé ajouter ?
  \item On tire ensuite un câble direct entre le PC et la télévision, ce qui crée un cycle
        \lstinline{Box-PC-TV}. La recherche
        \lstinline{peuvent_communiquer(reseau, "TV", "Tablette")} se termine-t-elle ? Que
        renvoie-t-elle ?
\end{enumerate}
\end{exercice}

% BEGIN-VERIFY
% def liens_non_reciproques(reseau):
%     defauts = []
%     for machine, voisins in reseau.items():
%         for voisin in voisins:
%             if machine not in reseau.get(voisin, []):
%                 defauts.append((machine, voisin))
%     return defauts
% def peuvent_communiquer(reseau, depart, arrivee, visites=None):
%     if visites is None:
%         visites = set()
%     if depart == arrivee:
%         return True
%     visites.add(depart)
%     for voisin in reseau.get(depart, []):
%         if voisin not in visites:
%             if peuvent_communiquer(reseau, voisin, arrivee, visites):
%                 return True
%     return False
% reseau = {"Box": ["PC", "TV", "Tablette"], "PC": ["Box"], "TV": []}
% assert peuvent_communiquer(reseau, "Box", "TV") is True
% assert peuvent_communiquer(reseau, "TV", "Box") is False
% assert liens_non_reciproques(reseau) == [("Box", "TV"), ("Box", "Tablette")]
% reseau["TV"] = ["Box"]
% reseau["Tablette"] = ["Box"]
% assert liens_non_reciproques(reseau) == []
% assert peuvent_communiquer(reseau, "TV", "Box") is True
% reseau["PC"].append("TV")
% reseau["TV"].append("PC")
% visites = set()
% assert peuvent_communiquer(reseau, "TV", "Tablette", visites) is True
% assert visites <= {"TV", "Box", "PC"}
% END-VERIFY
